\chapter{Variance, Moments, and Spread}

In the last chapter you learned to find the \emph{center} of a random
variable: its expectation, the long-run average value, the balance point of
its distribution. But two random quantities can share the same center and still
behave differently. A game that wins or loses a dollar with equal probability
has the same average payoff as one that wins or loses a million dollars. The
average is the same; the risk is not.

This chapter is about \emph{spread}: how far a random variable typically
wanders from its center. That spread is measured by \emph{variance}, and its
square root, the \emph{standard deviation}. We will see why variances of
independent variables add, how higher moments capture asymmetry and tail
weight, and how Jensen's inequality relates functions of averages to averages of
functions.

\chapterline

\section{The Problem: Two Distributions, One Center}
\label{ch09:problem}

Let us make the opening story concrete. Consider two games, each costing
nothing to play, each with expected payoff zero.

\begin{itemize}
  \item \textbf{Game A (mild).} You win \$1 or lose \$1, each with
        probability $1/2$.
  \item \textbf{Game B (wild).} You win \$100 or lose \$100, each with
        probability $1/2$.
\end{itemize}

Call the payoffs $X$ and $Y$. Both have expectation zero:
\[
  \E[X] = \tfrac12(1) + \tfrac12(-1) = 0,
  \qquad
  \E[Y] = \tfrac12(100) + \tfrac12(-100) = 0.
\]
A summary that only reports the mean would call these games identical.
Yet nobody would treat these two games the same: Game B carries far greater risk.
We need a number that grows when the outcomes are flung far from the center and
stays small when they huddle close.

\begin{figure}[H]
\centering
\begin{tikzpicture}
\begin{axis}[
  name=plotA,
  width=6.3cm, height=5.0cm,
  ybar, bar width=12pt,
  ymin=0, ymax=0.65,
  xmin=-3, xmax=3,
  title style={font=\small},
  title={Game A (mild): $\mu=0, \sigma=1$},
  xlabel={payoff (\$)}, ylabel={probability},
  xtick={-2,-1,0,1,2},
  ytick={0,0.25,0.5},
]
\addplot[draw=sourceblue, fill=sourceblue!55] coordinates {(-1,0.5) (1,0.5)};
\node[anchor=south, font=\scriptsize] at (axis cs:-1,0.5) {$0.5$};
\node[anchor=south, font=\scriptsize] at (axis cs:1,0.5) {$0.5$};
\draw[dashed, sourceblue!80, thick] (axis cs:0,0) -- (axis cs:0,0.58)
  node[above, font=\scriptsize] {$\mu = 0$};
\end{axis}

\begin{axis}[
  at={($(plotA.east)+(1.0cm,0)$)},
  anchor=west,
  width=6.3cm, height=5.0cm,
  ybar, bar width=12pt,
  ymin=0, ymax=0.65,
  xmin=-150, xmax=150,
  title style={font=\small},
  title={Game B (wild): $\mu=0, \sigma=100$},
  xlabel={payoff (\$)}, ylabel={},
  xtick={-100,-50,0,50,100},
  ytick={0,0.25,0.5},
  yticklabels={},
]
\addplot[draw=densityred, fill=densityred!45] coordinates {(-100,0.5) (100,0.5)};
\node[anchor=south, font=\scriptsize] at (axis cs:-100,0.5) {$0.5$};
\node[anchor=south, font=\scriptsize] at (axis cs:100,0.5) {$0.5$};
\draw[dashed, densityred!80, thick] (axis cs:0,0) -- (axis cs:0,0.58)
  node[above, font=\scriptsize] {$\mu = 0$};
\end{axis}
\end{tikzpicture}
\caption{Two distributions with identical mean ($\mu = 0$) but vastly different
spread. Game A (left, blue) hugs the center within $\pm \$1$, carrying standard
deviation $\sigma = 1$. Game B (right, red) flings the same probability far out
to $\pm \$100$, producing $\sigma = 100$. A summary reporting only the mean cannot
distinguish them; variance and standard deviation capture the difference.}
\label{ch09:fig-twogames}
\end{figure}

\begin{intuitionbox}
Think of a distribution as a pile of sand sitting on a seesaw. The
expectation is the point where the seesaw balances. The \emph{spread} is how
spread out the sand is around that balance point. You could pile it all in a
narrow heap at the center, or smear it thin out to the edges --- the balance
point can be identical in both cases. We want a single number that measures
``how smeared out.''
\end{intuitionbox}

\subsection{Why not just average the distance?}

The natural first idea is: measure each outcome's distance from the mean, and
average those distances. If $\mu = \E[X]$, look at $\abs{X - \mu}$ and take
its expectation. This quantity, $\E\,\abs{X-\mu}$, is called the \emph{mean
absolute deviation}, and it is a perfectly sensible measure of spread. For
Game A it equals $1$; for Game B it equals $100$. So far so good.

But the absolute value has a sharp corner at zero, and corners are the enemy
of calculus. The function $\abs{x}$ is not differentiable at the origin,
which makes the mean absolute deviation awkward to optimize, awkward to
combine across several random variables, and awkward to push through the
machinery of sums and limits we will build in later chapters. Mathematicians
and statisticians, faced with this corner, made a fateful and enormously
productive choice: instead of penalizing distance by $\abs{x-\mu}$, penalize
it by $(x-\mu)^2$. Squaring kills the corner --- it is smooth everywhere ---
and, as we will see, it interacts cleanly with sums, with independence,
and with the geometry of the normal distribution. The price is that squaring
exaggerates large deviations and gives answers in squared units. We pay that
price gladly, and then take a square root at the end to get back to the
original units.

\begin{recoverybox}
\textbf{Pause and think:} Suppose a coin flip pays $+\$5$ on heads and $-\$5$ on tails.
What is its mean absolute deviation? Its center is $\mu = 0$, and every outcome is exactly
$5$ units away, so $\E[\abs{X-\mu}] = 5$. What is its variance?
The squared deviation is $(5)^2 = 25$ every time, so $\Var(X) = 25$ dollars-squared,
giving standard deviation $\SD(X) = \sqrt{25} = 5$ dollars. Notice how the square root
faithfully restores the original unit and scale of fluctuation.
\end{recoverybox}

\section{Variance and Standard Deviation}
\label{ch09:defn}

\begin{definition}[Variance]
\label{ch09:def-var}
Let $X$ be a random variable with finite mean $\mu = \E[X]$. The
\emph{variance} of $X$ is the expected squared deviation from the mean,
\[
  \Var(X) \;=\; \E\!\big[(X-\mu)^2\big].
\]
Its non-negative square root is the \emph{standard deviation},
\[
  \SD(X) \;=\; \sqrt{\Var(X)},
\]
often written $\sigma$, with $\sigma^2$ for the variance.
\end{definition}

Three properties follow at once.

First, \textbf{variance is never negative.} It is the average of squares, and
squares are never negative, so $\Var(X) \ge 0$ always. It equals zero
exactly when $(X-\mu)^2 = 0$ with probability one --- that is, when $X$ is a
constant. A constant has no spread, and variance correctly reports zero. This
is your first sanity check on any variance calculation: a negative answer is
always an arithmetic mistake.

Second, \textbf{variance lives in squared units.} If $X$ is measured in
dollars, $\Var(X)$ is in dollars-squared, which is meaningless to a human. The
standard deviation $\SD(X)$ undoes the squaring and returns to dollars. That
is why, when reporting results to people, we almost always quote the standard
deviation, not the variance --- ``the returns have a standard deviation of
\$3'' is something you can feel, while ``a variance of 9 dollars-squared'' is
not.

Third, \textbf{the standard deviation is a genuine yardstick.} It tells you
the typical scale of fluctuation. A standard deviation of $\$3$ means that
draws of $X$ routinely land a few dollars away from the mean; landing
$\$30$ away would be extraordinary. We will sharpen ``routinely'' and
``extraordinary'' into a precise statement in Section~\ref{ch09:chebyshev}
(Chebyshev's inequality) but the slogan to carry is: \emph{the standard
deviation is the natural unit of surprise.}

\begin{example}[The two games, measured]
\label{ch09:ex-twogames}
Return to Games A and B from Section~\ref{ch09:problem}, both with mean
$\mu=0$.
For Game A,
\[
  \Var(X) = \tfrac12(1-0)^2 + \tfrac12(-1-0)^2 = \tfrac12 + \tfrac12 = 1,
  \qquad \SD(X) = 1.
\]
For Game B,
\[
  \Var(Y) = \tfrac12(100)^2 + \tfrac12(-100)^2 = 10{,}000,
  \qquad \SD(Y) = 100.
\]
The standard deviations are $1$ and $100$: Game B has one hundred times
the spread, matching the hundredfold jump in stakes. Notice that the
\emph{variances} differ by a factor of $10{,}000$, the square of the scale
factor. Squaring magnifies; the square root in the standard deviation brings
it back to a linear scale matching the original units.
\end{example}

\begin{intuitionbox}
\textbf{The mechanical picture: variance as moment of inertia.}
Why penalize distance with the square $(x-\mu)^2$ rather than another power?
In mechanics, placing weights $m_i$ along a rod locates the center of gravity at
$\mu = \sum m_i x_i$ (the expectation). Now pinch the rod at that pivot $\mu$ and try
to spin it! How much does the system resist angular acceleration?

The answer from physics is the \emph{moment of inertia}:
\[
  I = \sum m_i (x_i - \mu)^2.
\]
Masses huddled near the pivot have tiny lever arms $r_i = \abs{x_i - \mu}$, offering
almost zero resistance to spinning. But masses flung far out to the periphery resist
rotation with a lever arm that enters \emph{squared}! Variance $\Var(X) = \E[(X-\mu)^2]$
is literally the moment of inertia of the unit probability mass rotating about its
center of gravity. A low-variance variable spins with the slightest flick; a
high-variance variable carries massive rotational inertia.
\end{intuitionbox}

\begin{figure}[H]
\centering
\begin{tikzpicture}[font=\small, >=stealth]
  % Subfigure A: Low Variance
  \begin{scope}[shift={(-3.8,0)}]
    \node[font=\footnotesize\bfseries\color{sourceblue!80!black}] at (0, 2.6)
      {Panel (a): Low Variance $\Var(X) = 1$};
    % Vertical rotation axis
    \draw[dashed, thick, gray!70] (0,-0.15) -- (0,1.9);
    \node[above, font=\scriptsize\color{gray!90!black}] at (0,1.9) {Pivot $\mu = 0$};
    
    % Horizontal rod
    \draw[fill=gray!15, draw=gray!70, thick, rounded corners=1pt] (-2.2,-0.12) rectangle (2.2,0.12);
    
    % Clustered masses near center
    \draw[fill=sourceblue!35, draw=sourceblue!80, thick, rounded corners=2pt] (-0.9,-0.38) rectangle (-0.3,0.38);
    \node[above=3pt, font=\scriptsize\bfseries\color{sourceblue!90!black}] at (-0.6,0.38) {$m=\frac12$};
    \draw[fill=sourceblue!35, draw=sourceblue!80, thick, rounded corners=2pt] (0.3,-0.38) rectangle (0.9,0.38);
    \node[above=3pt, font=\scriptsize\bfseries\color{sourceblue!90!black}] at (0.6,0.38) {$m=\frac12$};
    
    % Lever arm dimension
    \draw[<->, sourceblue!80, semithick] (-0.6,-0.6) -- (0,-0.6);
    \draw[<->, sourceblue!80, semithick] (0,-0.6) -- (0.6,-0.6) node[midway, below=1pt, font=\scriptsize] {$r=1$};
    
    % Rotation indicator (easy spin)
    \draw[->, thick, targetgreen!80!black] (0.4, 1.25) arc (20:340:0.45cm and 0.15cm);
    \node[right=2pt, font=\scriptsize\color{targetgreen!80!black}] at (0.5, 1.25) {Spins easily};
    
    \node[draw=sourceblue!80, fill=sourceblue!10, rounded corners=2pt, font=\scriptsize, align=center] at (0,-1.45)
      {Moment of Inertia:\\[2pt]$I = \sum m_i r_i^2 = 1.0$};
  \end{scope}

  % Subfigure B: High Variance
  \begin{scope}[shift={(3.8,0)}]
    \node[font=\footnotesize\bfseries\color{densityred!80!black}] at (0, 2.6)
      {Panel (b): High Variance $\Var(Y) = 16$};
    % Vertical rotation axis
    \draw[dashed, thick, gray!70] (0,-0.15) -- (0,1.9);
    \node[above, font=\scriptsize\color{gray!90!black}] at (0,1.9) {Pivot $\mu = 0$};
    
    % Horizontal rod
    \draw[fill=gray!15, draw=gray!70, thick, rounded corners=1pt] (-2.6,-0.12) rectangle (2.6,0.12);
    
    % Distant masses flung to edges
    \draw[fill=densityred!35, draw=densityred!80, thick, rounded corners=2pt] (-2.5,-0.38) rectangle (-1.9,0.38);
    \node[above=3pt, font=\scriptsize\bfseries\color{densityred!90!black}] at (-2.2,0.38) {$m=\frac12$};
    \draw[fill=densityred!35, draw=densityred!80, thick, rounded corners=2pt] (1.9,-0.38) rectangle (2.5,0.38);
    \node[above=3pt, font=\scriptsize\bfseries\color{densityred!90!black}] at (2.2,0.38) {$m=\frac12$};
    
    % Lever arm dimension
    \draw[<->, densityred!80, semithick] (-2.2,-0.6) -- (-0.05,-0.6) node[midway, below=1pt, font=\scriptsize] {$r=4$};
    \draw[<->, densityred!80, semithick] (0.05,-0.6) -- (2.2,-0.6) node[midway, below=1pt, font=\scriptsize] {$r=4$};
    
    % Rotation indicator (high resistance)
    \draw[->, very thick, trapred] (0.5, 1.25) arc (20:160:0.55cm and 0.2cm);
    \node[right=2pt, font=\scriptsize\color{trapred!90!black}] at (0.6, 1.25) {High resistance};
    
    \node[draw=densityred!80, fill=densityred!10, rounded corners=2pt, font=\scriptsize, align=center] at (0,-1.45)
      {Moment of Inertia:\\[2pt]$I = \sum m_i r_i^2 = 16.0$};
  \end{scope}
\end{tikzpicture}
\caption{Variance as moment of inertia. Just as expectation $\mu = \E[X]$ is the center of gravity along a balance beam, variance $\Var(X) = \E[(X-\mu)^2]$ is the mechanical moment of inertia resisting rotation about that center. Concentrated mass (left) spins with negligible torque; mass flung out to the periphery (right) strongly resists angular acceleration because lever arms enter squared.}
\label{ch09:fig-moment-of-inertia}
\end{figure}

\begin{example}[A single fair die]
\label{ch09:ex-die}
Let $X$ be the result of one roll of a fair six-sided die. Its mean is
$\mu = (1+2+3+4+5+6)/6 = 3.5$. The variance is the average of the squared
deviations:
\[
  \Var(X) = \frac{1}{6}\sum_{k=1}^{6}(k-3.5)^2
  = \frac{1}{6}\big(6.25+2.25+0.25+0.25+2.25+6.25\big)
  = \frac{17.5}{6} \approx 2.92.
\]
So $\SD(X) = \sqrt{2.9167} \approx 1.71$. A typical roll lands about $1.7$
away from $3.5$, which matches the feel of the numbers $1$ through $6$ spread
evenly around $3.5$.
\end{example}

\section{The Computational Formula}
\label{ch09:compformula}

Computing variance straight from the definition means first finding $\mu$,
then forming $(X-\mu)^2$ for every outcome, then averaging. That is two passes
over the distribution. There is a slicker route that needs only the first two
\emph{moments} --- $\E[X]$ and $\E[X^2]$ --- and it is the formula you will
use ninety percent of the time.

\begin{theorem}[Computational formula for variance]
\label{ch09:thm-comp}
For any random variable $X$ with finite second moment,
\[
  \Var(X) \;=\; \E[X^2] - \big(\E[X]\big)^2.
\]
\end{theorem}

\begin{proof}
Expand the square inside the definition and use linearity of expectation
(Chapter ``Expectation''). Writing $\mu = \E[X]$, which is a constant,
\[
  \Var(X) = \E\big[(X-\mu)^2\big]
  = \E\big[X^2 - 2\mu X + \mu^2\big].
\]
Linearity lets us split the expectation over the sum, and pull the constants
$2\mu$ and $\mu^2$ outside:
\[
  = \E[X^2] - 2\mu\,\E[X] + \mu^2
  = \E[X^2] - 2\mu\cdot\mu + \mu^2
  = \E[X^2] - \mu^2.
\]
That is exactly $\E[X^2] - (\E[X])^2$. \qedhere
\end{proof}

\begin{whybox}
Why does ``mean of the square minus square of the mean'' measure spread? Read
it as a competition between two ways of averaging. $\E[X^2]$ is the average of
the squares; $(\E[X])^2$ is the square of the average. These two are equal
only when $X$ never varies. The moment $X$ spreads out, the big values get
\emph{extra} weight under squaring (because squaring is convex --- more on
this when we meet Jensen's inequality), so $\E[X^2]$ pulls ahead of
$(\E[X])^2$. The gap between them is exactly the spread: variance measures how
much squaring magnifies the dispersion of $X$ away from its mean.
\end{whybox}

\begin{trapbox}
A perennial error is to compute $(\E[X])^2$ and call it $\E[X^2]$. These are
almost never equal! For the fair die, $\E[X]=3.5$ so $(\E[X])^2 = 12.25$,
while $\E[X^2] = \tfrac16(1+4+9+16+25+36) = \tfrac{91}{6} \approx 15.17$. The
difference $15.17 - 12.25 = 2.92$ is the variance --- matching
Example~\ref{ch09:ex-die}. If you ever find $\E[X^2] < (\E[X])^2$, you have
made an arithmetic slip, because variance cannot be negative.
\end{trapbox}

\begin{example}[Bernoulli variance]
\label{ch09:ex-bern}
Let $X \sim \Bern(p)$, so $X=1$ with probability $p$ and $X=0$ with
probability $1-p$. Then $\E[X] = p$. Because $X$ only takes the values $0$ and
$1$, we have $X^2 = X$, so $\E[X^2] = \E[X] = p$ as well. The computational
formula gives
\[
  \Var(X) = \E[X^2] - (\E[X])^2 = p - p^2 = p(1-p).
\]
This little formula $p(1-p)$ is one of the most important in all of
statistics. Plotted against $p$ it is a downward parabola, zero at $p=0$ and
$p=1$ (a coin that always lands the same way has no spread) and largest at
$p=1/2$, where uncertainty is maximal.
\end{example}

\begin{figure}[H]
\centering
\begin{tikzpicture}
\begin{axis}[
  width=11cm, height=5.2cm,
  domain=0:1, samples=120,
  title style={font=\small},
  title={Variance of a Bernoulli trial is largest at maximum uncertainty},
  xlabel={$p$}, ylabel={$\Var(X)=p(1-p)$},
  ymin=0, ymax=0.35,
  xtick={0,0.25,0.5,0.75,1},
  ytick={0,0.125,0.25},
  grid=both, grid style={gray!18},
]
\addplot[densityred, very thick] {x*(1-x)};
\addplot[only marks, mark=*, sourceblue] coordinates {(0.5,0.25)};
\node[anchor=south, fill=white, inner sep=2pt] at (axis cs:0.5,0.25) {\small $\left(\tfrac12,\tfrac14\right)$};
\end{axis}
\end{tikzpicture}
\caption{The Bernoulli variance $p(1-p)$. A fair coin ($p=1/2$) is the most
unpredictable and carries the most variance; a coin guaranteed to land one
way ($p=0$ or $1$) carries none.}
\label{ch09:fig-bern}
\end{figure}

\begin{recoverybox}
\textbf{Pause and think:} Look closely at the Bernoulli variance formula $\Var(X) = p(1-p)$.
Why does it vanish at $p=0$ and $p=1$? If a coin has two heads ($p=1$), every toss is
guaranteed to show heads with probability $1$. There is no surprise and no variation,
so the spread is zero. Maximum spread occurs at $p=1/2$, where uncertainty is maximal,
giving $\Var(X) = 1/4$ and standard deviation $\SD(X) = 1/2$.
\end{recoverybox}

\section{How Variance Behaves Under Shifts and Scalings}
\label{ch09:affine}

To work fluently with variance you need to know what happens when you slide a
random variable along the number line or stretch it. These two rules come up
constantly and follow directly from the definition.

\begin{proposition}[Affine transformation]
\label{ch09:prop-affine}
For any constants $a$ and $b$,
\[
  \Var(aX + b) = a^2 \Var(X),
  \qquad\text{equivalently}\qquad
  \SD(aX+b) = \abs{a}\,\SD(X).
\]
\end{proposition}

\begin{proof}
Let $\mu = \E[X]$, so $\E[aX+b] = a\mu + b$ by linearity. Then
\[
  \Var(aX+b) = \E\Big[\big((aX+b) - (a\mu+b)\big)^2\Big]
  = \E\Big[\big(a(X-\mu)\big)^2\Big]
  = a^2\,\E\big[(X-\mu)^2\big] = a^2\Var(X). \qedhere
\]
\end{proof}

Two messages live in this one line.

\textbf{Shifting does nothing.} The constant $b$ vanished. Adding a fixed
amount to every outcome slides the whole distribution rigidly along the axis
without changing how spread out it is. If everyone in a room gets a \$10
raise, the inequality between their salaries is untouched. Variance, being a
measure of spread, ignores the location entirely.

\textbf{Scaling squares.} Multiplying by $a$ multiplies the standard
deviation by $\abs{a}$ and the variance by $a^2$. Double every outcome and you
double the spread (standard deviation $\times 2$) but quadruple the variance.
This is the same factor-of-square we saw between Games A and B: stretching the
axis by $100$ multiplied the variance by $10{,}000$.

\begin{happeningbox}
The appearance of $a^2$ (not $a$) is the clearest sign that variance is a
\emph{quadratic} quantity --- it behaves like an area, not a length. Standard
deviation is the corresponding length. This is exactly the relationship
between a square's area and its side. When you later see variances adding like
the squares of legs in the Pythagorean theorem
(Section~\ref{ch09:sumindep}), remember that variance was a squared quantity
from birth.
\end{happeningbox}

\section{Standardization and the z-Score}
\label{ch09:zscore}

Proposition~\ref{ch09:prop-affine} hands us a useful construction. Given any
random variable $X$ with mean $\mu$ and standard deviation $\sigma > 0$, we
can manufacture a new variable that is centered at zero and has spread
exactly one. Subtract the mean (to recenter) and divide by the standard
deviation (to rescale):
\[
  Z \;=\; \frac{X - \mu}{\sigma}.
\]
This is the \emph{standardization} of $X$, and $Z$ is called the
\emph{z-score}. Let us verify its two defining properties. By linearity,
\[
  \E[Z] = \frac{\E[X] - \mu}{\sigma} = \frac{\mu - \mu}{\sigma} = 0,
\]
and by the affine rule with $a = 1/\sigma$,
\[
  \Var(Z) = \frac{1}{\sigma^2}\Var(X) = \frac{\sigma^2}{\sigma^2} = 1.
\]
So every $X$, no matter its original units --- dollars, kilograms, exam
points --- can be converted into a dimensionless number $Z$ whose value is
``how many standard deviations above or below the mean am I?''

\begin{intuitionbox}
A z-score is a universal translator for surprise. A student who scores $650$
on one exam and $34$ on another cannot be compared until you know the mean and
spread of each test. But if both scores correspond to $z = +1.5$, you know
instantly that the student did equally well on both relative to peers ---
one-and-a-half standard deviations above average in each case. The z-score
strips away the arbitrary units and leaves pure standing.
\end{intuitionbox}

\begin{example}[Two exams]
\label{ch09:ex-exams}
On a midterm with mean $70$ and standard deviation $8$, you score $82$. On a
final with mean $65$ and standard deviation $20$, you score $90$. Which is the
better performance relative to the class?
\[
  z_{\text{mid}} = \frac{82-70}{8} = 1.5,
  \qquad
  z_{\text{fin}} = \frac{90-65}{20} = 1.25.
\]
The raw final score ($90$) is higher, but the midterm z-score ($1.5$) is
larger: you stood out more on the midterm, because that test had much less
spread, so $12$ points above its mean was rarer than $25$ points above the
final's mean.
\end{example}

\begin{figure}[H]
\centering
\begin{tikzpicture}[font=\small, >=stealth]
  % Top axis: raw scale X
  \draw[->, thick, gray!80] (3.8, 2.2) -- (12.2, 2.2) node[right, font=\footnotesize] {raw score $X$};
  \foreach \x/\val in {4.6/{$54$}, 5.8/{$62$}, 7.0/{$70$}, 8.2/{$78$}, 9.4/{$86$}, 10.6/{$94$}} {
    \draw[thick, gray!80] (\x, 2.05) -- (\x, 2.35);
    \node[below=2pt, font=\scriptsize] at (\x, 2.05) {\val};
  }
  \node[above=2pt, font=\scriptsize\bfseries\color{sourceblue!80!black}] at (7.0, 2.35) {$\mu = 70$};
  \draw[<->, thick, sourceblue!70] (7.0, 2.95) -- (8.2, 2.95);
  \node[above=2pt, font=\tiny] at (7.6, 2.95) {$\sigma = 8$};

  % Highlight raw point X = 82: 7.0 + 1.5 * 1.2 = 8.8
  \filldraw[sourceblue] (8.8, 2.2) circle (2.5pt);
  \node[above=3pt, font=\footnotesize\bfseries\color{sourceblue}] at (8.8, 2.2) {$X = 82$};

  % Transformation arrows in between
  \draw[->, thick, dashed, quotientorange!90!black] (7.0, 1.7) -- (7.0, 0.95)
    node[midway, left=2pt, font=\scriptsize] {Shift $-\mu$};
  \draw[->, thick, dashed, quotientorange!90!black] (8.8, 1.7) -- (8.8, 0.95)
    node[midway, right=2pt, font=\scriptsize] {Rescale $\div\sigma$};

  % Bottom axis: standardized scale Z
  \draw[->, thick, gray!80] (3.8, 0.2) -- (12.2, 0.2) node[right, font=\footnotesize] {z-score $Z$};
  \foreach \x/\val in {4.6/{$-2$}, 5.8/{$-1$}, 7.0/{$0$}, 8.2/{$+1$}, 9.4/{$+2$}, 10.6/{$+3$}} {
    \draw[thick, gray!80] (\x, 0.05) -- (\x, 0.35);
    \node[below=2pt, font=\scriptsize] at (\x, 0.05) {\val};
  }
  \node[below=10pt, font=\scriptsize\bfseries\color{targetgreen!80!black}] at (7.0, 0.05) {Mean $= 0$};
  \draw[<->, thick, targetgreen!80!black] (7.0, -0.75) -- (8.2, -0.75);
  \node[below=2pt, font=\tiny] at (7.6, -0.75) {$\SD = 1$};

  % Highlight standardized point Z = +1.5
  \filldraw[densityred] (8.8, 0.2) circle (2.5pt);
  \node[above=3pt, font=\footnotesize\bfseries\color{densityred}] at (8.8, 0.2) {$Z = +1.5$};
\end{tikzpicture}
\caption{Standardization transforms any measurement scale into standard units.
Subtracting the mean $\mu = 70$ slides the distribution so its center lands at $0$.
Dividing by the standard deviation $\sigma = 8$ rescales the axis so that one unit
of distance corresponds exactly to one standard deviation. A raw midterm score of $82$
becomes a z-score of $z = (82-70)/8 = +1.5$, visibly placing it $1.5$ standard
deviations above average.}
\label{ch09:fig-zscore}
\end{figure}

Figure~\ref{ch09:fig-zscore} illustrates this transformation: subtracting
$70$ slides the center from $70$ to zero, and dividing by $8$ recalibrates
the tick marks so that every unit of distance represents one full standard deviation.
Standardization is not a mere convenience; it is the gateway to the normal
distribution and the central limit theorem in later chapters. There, the
standardized version of a sum of many independent pieces converges to a fixed
universal shape --- the standard normal --- regardless of the specific distributions
of the summands. Z-scores provide the natural coordinates in which this convergence
becomes apparent.

\section{Chebyshev's Inequality: Making ``Surprise'' Precise}
\label{ch09:chebyshev}

We have twice promised to turn the slogan ``the standard deviation is the unit
of surprise'' into a theorem. Here it is. The remarkable thing is that a single
inequality, demanding \emph{nothing} about the shape of the distribution ---
not symmetry, not unimodality (having only a single central peak), not even a
known formula --- already forces most of the probability mass to live within a
few standard deviations of the mean.

\begin{theorem}[Chebyshev's inequality]
\label{ch09:thm-cheby}
Let $X$ have finite mean $\mu$ and finite variance $\sigma^2$. For any
$k > 0$,
\[
  \P\big(\abs{X-\mu} \ge k\sigma\big) \;\le\; \frac{1}{k^2}.
\]
In words: the probability of landing $k$ or more standard deviations away from
the mean is at most $1/k^2$.
\end{theorem}

\begin{proof}
Let $D = (X-\mu)^2$. Notice that $D$ is a non-negative random variable ($D \ge 0$)
whose expectation is the variance: $\E[D] = \sigma^2$. The event that $X$ wanders
$k$ standard deviations from the mean,
\[
  \{\abs{X-\mu} \ge k\sigma\},
\]
is identical to the event that its squared deviation is at least $k^2\sigma^2$:
\[
  \{D \ge k^2\sigma^2\}.
\]
Now we use the elementary step-function bound known as \emph{Markov's inequality}.
For any non-negative random variable $W \ge 0$ and any positive threshold $t > 0$,
the indicator $\ind\{W \ge t\}$ (which is $1$ when $W \ge t$ and $0$ otherwise)
satisfies
\[
  \ind\{W \ge t\} \;\le\; \frac{W}{t}.
\]
Check this by cases: when $W < t$, the indicator is $0$ while $W/t \ge 0$; when
$W \ge t$, the indicator is $1$ while $W/t \ge 1$. Taking expectations of both
sides yields Markov's inequality: $\P(W \ge t) \le \E[W]/t$.

Applying this with $W = D$ and threshold $t = k^2\sigma^2$, we obtain
\[
  \P\big(\abs{X-\mu} \ge k\sigma\big)
  \;=\; \P\big(D \ge k^2\sigma^2\big)
  \;=\; \E\big[\ind\{D \ge k^2\sigma^2\}\big]
  \;\le\; \frac{\E[D]}{k^2\sigma^2}
  \;=\; \frac{\sigma^2}{k^2\sigma^2}
  \;=\; \frac{1}{k^2}. \qedhere
\]
\end{proof}

Setting $k=2$ shows at least $1 - 1/4 = 75\%$ of any distribution's mass lies
within two standard deviations of the mean; $k=3$ guarantees at least
$1 - 1/9 \approx 89\%$ within three. These bounds are deliberately
\emph{loose} --- they must hold for every distribution at once, including
adversarial ones engineered to push mass outward --- so for any particular
well-behaved distribution the truth is far tighter.

\begin{example}[Chebyshev versus the normal rule]
\label{ch09:ex-cheby-vs-normal}
For the bell-shaped normal distribution the famous ``empirical rule'' says
about $68\%$ of mass lies within one standard deviation, $95\%$ within two, and
$99.7\%$ within three. Compare these to what Chebyshev guarantees for
\emph{any} distribution:
\[
  \begin{array}{c|c|c}
    k & \text{Chebyshev bound on } \P(\abs{X-\mu}<k\sigma)
      & \text{normal (empirical rule)}\\\hline
    2 & \ge 75\%   & \approx 95\%\\
    3 & \ge 88.9\% & \approx 99.7\%
  \end{array}
\]
The normal does far better than Chebyshev demands --- that is the reward for
having thin, light tails. Chebyshev is the universal floor; the normal rule is
what a particularly tame distribution actually achieves.
\end{example}

\begin{figure}[H]
\centering
\begin{tikzpicture}
\begin{axis}[
  width=12.5cm, height=5.4cm,
  domain=-4:4, samples=160,
  title style={font=\small},
  title={The normal density with its tails beyond $\pm 2\sigma$ shaded},
  xlabel={standard deviations from the mean ($z$)},
  ylabel={density},
  ymin=0, ymax=0.45,
  xtick={-3,-2,-1,0,1,2,3},
  xticklabels={$\mu-3\sigma$,$\mu-2\sigma$,$\mu-\sigma$,$\mu$,$\mu+\sigma$,$\mu+2\sigma$,$\mu+3\sigma$},
  ytick=\empty,
]
\addplot[name path=curve, distblue, very thick, domain=-4:4]
  {0.3989*exp(-x^2/2)};
\path[name path=axis] (axis cs:-4,0) -- (axis cs:4,0);
\addplot[densityred!35] fill between[of=curve and axis, soft clip={domain=-4:-2}];
\addplot[densityred!35] fill between[of=curve and axis, soft clip={domain=2:4}];
\node at (axis cs:2.75,0.05) {\small tail};
\node at (axis cs:-2.75,0.05) {\small tail};
\node at (axis cs:0,0.2) {\small $\approx 95\%$};
\end{axis}
\end{tikzpicture}
\caption{The standard normal density. The shaded tails beyond $\pm 2\sigma$
together hold only about $5\%$ of the mass, so roughly $95\%$ sits within two
standard deviations. Chebyshev's inequality guarantees \emph{at least} $75\%$
for any distribution whatsoever; the bell curve, with its light tails,
comfortably beats that floor. The standard deviation is the ruler along the
horizontal axis.}
\label{ch09:fig-bell}
\end{figure}

\begin{aestheticbox}
Chebyshev is the reason variance is more than bookkeeping. A small variance
does not merely \emph{describe} a tidy distribution --- it \emph{forces} one.
Once you know $\sigma$, you have a guaranteed, distribution-free leash on how
far the variable can roam: the probability of straying $k$ standard deviations
out can never exceed $1/k^2$, no matter how exotic the underlying randomness.
This is what lets you attach reliable error bars to estimates whose exact
distribution you do not know, and it is the inequality that powers the first
proof of the law of large numbers in a later chapter.
\end{aestheticbox}

\section{The Variance of a Sum and the Role of Independence}
\label{ch09:sumindep}

Expectation is remarkably simple under addition: $\E[X+Y] = \E[X] + \E[Y]$,
holding without any condition on independence. Variance, by contrast, behaves
differently under sums --- and the extra term that appears reveals how the variables interact.

\begin{theorem}[Variance of a sum of two variables]
\label{ch09:thm-sum2}
For any two random variables $X$ and $Y$ with finite variances,
\[
  \Var(X+Y) = \Var(X) + \Var(Y) + 2\,\Cov(X,Y),
\]
where the \emph{covariance} is
$\Cov(X,Y) = \E\big[(X-\mu_X)(Y-\mu_Y)\big]
= \E[XY] - \E[X]\E[Y]$.
\end{theorem}

\begin{proof}
Let $\mu_X = \E[X]$, $\mu_Y = \E[Y]$, so $\E[X+Y] = \mu_X + \mu_Y$ by linearity.
Abbreviating the centered variables by $\tilde X = X-\mu_X$ and
$\tilde Y = Y-\mu_Y$, the variance of the sum is
\[
  \Var(X+Y) = \E\big[(\tilde X + \tilde Y)^2\big]
  = \E\big[\tilde X^2 + 2\tilde X\tilde Y + \tilde Y^2\big].
\]
By linearity this splits into
\[
  \E[\tilde X^2] + 2\,\E[\tilde X \tilde Y] + \E[\tilde Y^2]
  = \Var(X) + 2\Cov(X,Y) + \Var(Y).
\]
Expanding the product inside the covariance gives the standard computational formula:
\[
  \Cov(X,Y) = \E\big[(X-\mu_X)(Y-\mu_Y)\big]
  = \E\big[XY - \mu_X Y - \mu_Y X + \mu_X\mu_Y\big]
  = \E[XY] - \E[X]\E[Y]. \qedhere
\]
\end{proof}

The cross term $2\Cov(X,Y)$ is the whole story. Covariance measures how $X$
and $Y$ move \emph{together}: it is positive when they tend to be large at
the same time, negative when one being large signals the other being small,
and zero when their fluctuations are unrelated in the averaged sense. We
devote a full later chapter to covariance and its normalized cousin
correlation; here we need only one consequence.

\begin{corollary}[Variances of independent variables add]
\label{ch09:cor-indepadd}
If $X$ and $Y$ are independent, then we have $\Cov(X,Y) = 0$, and therefore
\[
  \Var(X+Y) = \Var(X) + \Var(Y).
\]
More generally, if $X_1, \dots, X_n$ are pairwise independent (or merely
pairwise uncorrelated), then
\[
  \Var\!\Big(\sum_{i=1}^n X_i\Big) = \sum_{i=1}^n \Var(X_i).
\]
\end{corollary}

\begin{proof}
If $X \indep Y$ then $\E[XY] = \E[X]\E[Y]$ (a fact proved in the chapter on
expectation), so $\Cov(X,Y) = \E[XY] - \E[X]\E[Y] = 0$. The cross terms in
the general expansion $\Var(\sum_i X_i) = \sum_i \Var(X_i)
+ \sum_{i\ne j}\Cov(X_i,X_j)$ all vanish under pairwise uncorrelatedness,
leaving only the diagonal sum. \qedhere
\end{proof}

\begin{focusbox}
\textbf{Independence is exactly the condition that lets variances add.} The
mean is always additive. Variances add \emph{only when the variables are
uncorrelated}. This key distinction between mean and variance is central to
statistical theory --- error bars, the $\sqrt{n}$ law, portfolio diversification,
and the central limit theorem all depend on it.
\end{focusbox}

\begin{connectionbox}
Theorem~\ref{ch09:thm-sum2} is the Pythagorean theorem in the geometry of random variables. Picture
each centered random variable as a vector, with the ``length-squared'' of a
vector being its variance and the ``dot product'' of two vectors being their
covariance. Independent (uncorrelated) variables are then \emph{perpendicular}
vectors, and for perpendicular vectors the squared length of the sum is the
sum of squared lengths --- exactly $\Var(X+Y) = \Var(X) + \Var(Y)$. The
companion linear-algebra reader develops this inner-product picture; here it
is enough to know that the geometry is real: standard deviations
combine like the legs of a right triangle, $\sigma_{X+Y} =
\sqrt{\sigma_X^2 + \sigma_Y^2}$, as illustrated in Figure~\ref{ch09:fig-pythagoras}.
\end{connectionbox}

\begin{figure}[H]
\centering
\begin{tikzpicture}[font=\small, >=stealth]
  % Subfigure A: Independent (Orthogonal)
  \begin{scope}[shift={(-3.8,0)}]
    \node[font=\footnotesize\bfseries\color{sourceblue!80!black}] at (1.5, 3.1)
      {Panel (a): Uncorrelated ($\Cov = 0$)};

    \coordinate (O) at (0,0);
    \coordinate (X) at (3.0,0);
    \coordinate (Y) at (0,2.0);
    \coordinate (S) at (3.0,2.0);

    % Right angle marker
    \draw[thick, gray!70] (0.35,0) -- (0.35,0.35) -- (0,0.35);

    % Vector X
    \draw[->, very thick, sourceblue] (O) -- (X)
      node[midway, below=2pt, font=\scriptsize] {$\tilde X$};
    % Vector Y
    \draw[->, very thick, sourceblue] (O) -- (Y)
      node[midway, left=2pt, font=\scriptsize] {$\tilde Y$};
    % Dashed parallelogram
    \draw[dashed, gray!60] (X) -- (S);
    \draw[dashed, gray!60] (Y) -- (S);
    % Sum vector
    \draw[->, very thick, densityred] (O) -- (S)
      node[midway, above left=-1pt, font=\scriptsize] {$\tilde X + \tilde Y$};

    \node[draw=densityred!80, fill=densityred!8, rounded corners=2pt, font=\scriptsize, align=center] at (1.5,-1.15)
      {Right triangle ($90^\circ$):\\$\Var(X+Y) = \Var(X) + \Var(Y)$};
  \end{scope}

  % Subfigure B: Correlated
  \begin{scope}[shift={(3.4,0)}]
    \node[font=\footnotesize\bfseries\color{quotientorange!80!black}] at (1.6, 3.1)
      {Panel (b): Correlated ($\Cov > 0$)};

    \coordinate (O2) at (0,0);
    \coordinate (X2) at (2.8,0);
    \coordinate (Y2) at (1.2,1.8);
    \coordinate (S2) at ($(X2)+(Y2)$);

    \draw[->, thick, sourceblue] (O2) -- (X2) node[below=2pt, font=\scriptsize] {$\tilde X$};
    \draw[->, thick, sourceblue] (O2) -- (Y2) node[above left=-2pt, font=\scriptsize] {$\tilde Y$};
    \draw[dashed, gray!60] (X2) -- (S2);
    \draw[dashed, gray!60] (Y2) -- (S2);
    \draw[->, very thick, quotientorange] (O2) -- (S2)
      node[midway, above left=-2pt, font=\scriptsize] {$\tilde X+\tilde Y$};

    \node[right, font=\tiny\color{targetgreen!80!black}, align=center] at (2.4, 0.5)
      {Acute angle: $\Cov > 0$\\Stretched sum};

    \node[draw=quotientorange!80, fill=quotientorange!8, rounded corners=2pt, font=\scriptsize, align=center] at (1.6,-1.15)
      {Law of Cosines:\\$\Var(X+Y) = \Var(X) + \Var(Y) + 2\Cov(X,Y)$};
  \end{scope}
\end{tikzpicture}
\caption{The geometry of variance addition. When random variables are centered
($\tilde X = X-\mu_X$), variance acts as squared Euclidean length and covariance
as the inner product. Panel~(a): If $X$ and $Y$ are independent (uncorrelated),
their vectors are perpendicular ($90^\circ$); the squared length of the sum is the sum
of squared lengths, exactly as in the Pythagorean theorem ($\sigma_{X+Y} = \sqrt{\sigma_X^2 + \sigma_Y^2}$).
Panel~(b): When correlated, the angle tilts: positive covariance forms an acute angle,
stretching the sum (magnifying risk), while negative covariance forms an obtuse angle,
shrinking the sum (hedging).}
\label{ch09:fig-pythagoras}
\end{figure}

\subsection{The \texorpdfstring{$\sqrt{n}$}{sqrt(n)} law and why averaging works}

Here is the payoff that justifies all the machinery. Let $X_1, \dots, X_n$ be
independent draws from the same distribution, each with mean $\mu$ and
variance $\sigma^2$ --- think of $n$ repetitions of the same noisy
measurement. Consider their sum $S_n = \sum_i X_i$ and their average
$\bar X_n = S_n / n$.

For the \textbf{sum}, variances add (Corollary~\ref{ch09:cor-indepadd}):
\[
  \Var(S_n) = n\sigma^2,
  \qquad
  \SD(S_n) = \sqrt{n}\,\sigma.
\]
For the \textbf{average}, apply the scaling rule with $a = 1/n$:
\[
  \Var(\bar X_n) = \frac{1}{n^2}\Var(S_n) = \frac{n\sigma^2}{n^2}
  = \frac{\sigma^2}{n},
  \qquad
  \SD(\bar X_n) = \frac{\sigma}{\sqrt{n}}.
\]

\begin{aestheticbox}
This last line, $\SD(\bar X_n) = \sigma/\sqrt{n}$, is one of the most
consequential formulas in all of science. It says that averaging $n$
independent measurements shrinks the spread of your estimate by a factor of
$\sqrt{n}$. Average $4$ readings and the noise halves; average $100$ and it
drops by a factor of $10$. This is \emph{why repeated measurement works}, why
larger samples give tighter conclusions, and why every poll and every clinical
trial reports a margin of error that shrinks with the square root of the
sample size. The mean of the average stays put at $\mu$ (no bias), while its
spread melts away as $1/\sqrt{n}$ --- so the average homes in on the truth.
\end{aestheticbox}

\begin{figure}[H]
\centering
\begin{tikzpicture}
\begin{axis}[
  width=11.5cm, height=5.4cm,
  domain=1:100, samples=100,
  title style={font=\small},
  title={Diminishing returns: noise falls like $1/\sqrt{n}$},
  xlabel={number of measurements $n$},
  ylabel={$\SD(\bar X_n) = \sigma/\sqrt{n}$},
  ymin=0, ymax=1.05,
  xtick={1,25,50,75,100},
  grid=both, grid style={gray!18},
]
\addplot[sourceblue, very thick] {1/sqrt(x)};
\addplot[only marks, mark=*, densityred] coordinates {(1,1) (4,0.5) (25,0.2) (100,0.1)};
\node[anchor=west] at (axis cs:4,0.5) {\small $n=4\Rightarrow \tfrac12$};
\node[anchor=west] at (axis cs:25,0.2) {\small $n=25\Rightarrow \tfrac15$};
\end{axis}
\end{tikzpicture}
\caption{The standard deviation of a sample average decays as $\sigma/\sqrt
n$. The curve is steep at first --- the first few extra measurements help a
lot --- then flattens, so cutting the error in half always costs four times
the data. This diminishing return is a permanent feature of estimation.}
\label{ch09:fig-sqrtn}
\end{figure}

\begin{trapbox}
The $\sqrt{n}$ scaling has a sober flip side. To halve your error you must
\emph{quadruple} your sample, and to cut it to a tenth you need a hundred times
the data. Precision is expensive, and it gets more expensive the more you
already have. Beware anyone who promises that ``just collecting more data''
will cheaply rescue a noisy study; past a point, the $1/\sqrt n$ curve in
Figure~\ref{ch09:fig-sqrtn} is nearly flat.
\end{trapbox}

\begin{example}[Variance of the binomial]
\label{ch09:ex-binom}
A binomial random variable $X \sim \Bin(n,p)$ counts the successes in $n$
independent Bernoulli$(p)$ trials, so $X = \sum_{i=1}^n B_i$ with each
$B_i \sim \Bern(p)$ independent. Each $B_i$ has variance $p(1-p)$
(Example~\ref{ch09:ex-bern}), and because the trials are independent the
variances simply add:
\[
  \Var(X) = \sum_{i=1}^n \Var(B_i) = n\,p(1-p).
\]
Deriving this straight from the binomial pmf would mean wrestling with sums of
$k^2\binom{n}{k}p^k(1-p)^{n-k}$. Decomposing into independent indicators makes
it a one-liner. This ``break a complicated variable into a sum of simple
independent pieces'' move is one of the sharpest tactics in the subject.
\end{example}

\begin{applicationbox}[Why diversification reduces risk]
A portfolio's value is a weighted sum of its assets' returns, and a
portfolio's \emph{risk} is the standard deviation of that sum. If you put all
your money in one stock, your risk is that stock's full standard deviation. If
instead you split it across many assets whose returns are roughly
independent, the cross-covariances stay small and the variances add as in
Corollary~\ref{ch09:cor-indepadd} --- but you have also \emph{divided} each
position, so each scaled variance shrinks. The net effect is the $\sqrt n$
law again: spreading a fixed sum across $n$ uncorrelated assets cuts the
portfolio's standard deviation toward $\sigma/\sqrt n$. This is the
mathematical core of Harry Markowitz's portfolio theory, and the reason ``do
not put all your eggs in one basket'' is not folk wisdom but a theorem. The
catch sits in the covariance term: in a crash, assets fall \emph{together},
covariances turn sharply positive, and the diversification benefit
evaporates exactly when you need it most.
\end{applicationbox}

\begin{paradoxbox}[The independence assumption that wasn't]
The variance-adds rule is so convenient that people apply it reflexively ---
and forget its one hypothesis. Suppose an insurer writes $10{,}000$ flood
policies and reasons: each policy has small variance, variances add, so the
total variance is $10{,}000$ times small, and the standard deviation of the
payout is only $\sqrt{10{,}000} = 100$ times small --- a comfortably tiny
\emph{relative} risk by the $\sqrt n$ argument. Then a single flood soaks
every policyholder at once. The losses were never independent; they shared a
common cause. With perfect positive correlation the standard deviations
\emph{add} instead of the variances, so the spread grows like $n$, not
$\sqrt n$, and the insurer is wiped out. The moral: before you let variances
add, ask hard whether the parts can fail together. The $2\Cov$ term in
Theorem~\ref{ch09:thm-sum2} is not a footnote --- it is sometimes the entire
risk.
\end{paradoxbox}

\section{Moments: A Family of Summaries}
\label{ch09:moments}

The mean and variance are the first two members of an infinite family. The
$k$-th \emph{moment} of $X$ is $\E[X^k]$, and the $k$-th \emph{central
moment} is $\E[(X-\mu)^k]$, the moment taken about the mean.

\begin{definition}[Moments]
\label{ch09:def-moments}
For a positive integer $k$, the \emph{$k$-th raw moment} of $X$ is
$m_k = \E[X^k]$, and the \emph{$k$-th central moment} is
$\mu_k = \E[(X-\mu)^k]$, where $\mu = \E[X]$.
\end{definition}

What each moment captures follows a clear pattern:

\begin{itemize}
  \item \textbf{First moment} $\E[X] = \mu$ --- the \emph{center}.
  \item \textbf{Second central moment} $\E[(X-\mu)^2] = \Var(X)$ --- the
        \emph{spread}.
  \item \textbf{Third central moment} measures \emph{asymmetry} (skewness):
        is the distribution lopsided?
  \item \textbf{Fourth central moment} measures \emph{tail weight}
        (kurtosis): how prone is it to extreme values?
\end{itemize}

Each higher moment refines the silhouette of the distribution. The first
locates it, the second sizes it, the third tilts it, the fourth fattens or
thins its tails. In principle, knowing all the moments often pins down the
distribution completely --- a theme we return to with moment generating
functions in a later chapter.

\subsection{Skewness: the third moment, made scale-free}

The raw third central moment $\E[(X-\mu)^3]$ has units of $X$ cubed, which
makes it impossible to compare across different variables. We standardize it
by dividing by $\sigma^3$, producing a pure number.

\begin{definition}[Skewness]
\label{ch09:def-skew}
The \emph{skewness} of $X$ is
\[
  \operatorname{Skew}(X) = \E\!\left[\left(\frac{X-\mu}{\sigma}\right)^{3}\right]
  = \frac{\E[(X-\mu)^3]}{\sigma^3}.
\]
\end{definition}

The cube preserves sign, and that is the point. When $X$ exceeds its mean the
cubed deviation is positive; when it falls short it is negative. Skewness asks
which side wins on average. A symmetric distribution contributes equal
positive and negative cubes that cancel, giving skewness zero. A distribution
with a long \emph{right} tail --- a few very large values --- has its big
positive cubes dominate, giving \emph{positive} skew. A long left tail gives
negative skew.

\begin{intuitionbox}
The cleanest mental image: where is the long tail pointing? \emph{Right}-skewed
means the tail stretches toward large values, dragging the mean to the right
of the bulk. Incomes are the classic case --- most people earn modestly while
a few earn enormously, so the average income sits well above the median.
Whenever the mean noticeably exceeds the median, suspect a long right tail and
positive skew.
\end{intuitionbox}

\begin{figure}[H]
\centering
\begin{tikzpicture}
\begin{axis}[
  width=12.5cm, height=5cm,
  domain=0:10, samples=160,
  title style={font=\small},
  title={Right-skewed (long tail to the right) vs.\ symmetric},
  xlabel={$x$}, ylabel={density},
  ymin=0, ymax=0.65,
  legend pos=north east,
  grid=both, grid style={gray!15},
]
% a roughly symmetric bump centered near 5
\addplot[sourceblue, very thick] {0.55*exp(-(x-5)^2/1.6)};
\addlegendentry{symmetric}
% a right-skewed bump: gamma-ish shape
\addplot[densityred, very thick] {0.36*(x^2)*exp(-1.4*x)};
\addlegendentry{right-skewed}
\end{axis}
\end{tikzpicture}
\caption{A symmetric density (blue) has zero skewness: equal mass on each
side of its peak. A right-skewed density (red) rises quickly then trails off
in a long right tail, pulling its mean rightward and giving positive
skewness.}
\label{ch09:fig-skew}
\end{figure}

\subsection{Kurtosis: the fourth moment and the weight of the tails}

The fourth standardized central moment measures how much of a distribution's
variance comes from rare, extreme deviations rather than ordinary
middle-sized ones.

\begin{definition}[Kurtosis]
\label{ch09:def-kurt}
The \emph{kurtosis} of $X$ is
\[
  \operatorname{Kurt}(X) = \E\!\left[\left(\frac{X-\mu}{\sigma}\right)^{4}\right]
  = \frac{\E[(X-\mu)^4]}{\sigma^4}.
\]
The \emph{excess kurtosis} is $\operatorname{Kurt}(X) - 3$, defined so that the
normal distribution has excess kurtosis zero.
\end{definition}

Because the deviations are raised to the fourth power, an outlier four
standard deviations out contributes $4^4 = 256$ times as much as a typical
one-standard-deviation wobble. Kurtosis is therefore a tail-sensitivity meter.
High kurtosis (``heavy tails'') means the distribution, while perhaps calm
most of the time, produces occasional violent excursions. Low kurtosis means
the outcomes stay tightly within a typical band.

\begin{trapbox}
Kurtosis is \emph{not} a measure of ``peakedness,'' despite a stubborn old
textbook tradition that says so. Two distributions can have nearly identical
central peaks and vastly different kurtosis purely because of what happens far
out in the tails. The fourth power is overwhelmingly dominated by extreme
values, so read kurtosis as a statement about rare large events, not about
the shape of the hump in the middle.
\end{trapbox}

\begin{applicationbox}[Heavy tails and the limits of the bell curve]
The normal distribution has excess kurtosis exactly zero, and for decades
financial models assumed asset returns were normal. Real markets disagree
loudly: daily stock returns are sharply \emph{leptokurtic} (having heavy tails
and positive excess kurtosis, from the Greek \emph{leptos} for slender),
meaning crashes and extreme spikes occur far more often than a normal bell
curve predicts. A ``six-sigma'' daily move should, under normality, happen roughly
once in several million years; in real markets such moves arrive every few
years. Underestimating kurtosis --- pretending the tails are thin when they
are fat --- has bankrupted funds and helped fuel financial crises. The fourth
moment exposes how badly thin-tailed Gaussian models underestimate extreme
events, which is why modern risk management relies on heavy-tailed distributions
instead.
\end{applicationbox}

\begin{historybox}
The vocabulary of moments was borrowed directly from classical mechanics. In
physics, placing weights along a line yields the center of mass (the first
moment, $\sum m_i x_i$) and the moment of inertia (the second moment, $\sum m_i
r_i^2$, which measures physical resistance to rotation). In the 1850s and
1860s, mathematicians realized that a probability distribution could be viewed
as a physical bar of unit mass smeared across the real line.

Pafnuty Chebyshev (1821--1894) was wrestling with a roadblock:
earlier proofs of the Law of Large Numbers (by Jacob Bernoulli and Sim\'eon
Denis Poisson) relied on restrictive assumptions about coin flips or explicit
formulas. Probabilists lacked a mathematical tool to bound how far an arbitrary,
unknown distribution could stray from its center. In 1853, the French
statistician Ir\'en\'ee-Jules Bienaym\'e had noticed that independent variances
add, but it was Chebyshev in 1867 who turned the second moment into
his inequality, $\P(\abs{X-\mu} \ge k\sigma) \le 1/k^2$, giving the
first distribution-free proof of the Weak Law of Large Numbers.

Around 1900, the British statistician Karl Pearson (1857--1936) expanded this
mechanical framework, coining the word \emph{kurtosis} (Greek for bulging or
curvature) and developing the ``method of moments'' to classify biological
data. Yet moments brought controversy. Pearson originally claimed that kurtosis
measured the ``peakedness'' of a distribution's central hump. That intuition
misled researchers for half a century until Irving Kaplansky and others
proved that kurtosis is almost entirely determined by the distant tails, not
the middle peak. Pearson also clashed in the 1920s with Ronald Fisher,
who proved that estimating parameters by matching higher moments is noisy and
statistically inefficient compared to maximum likelihood (finding parameters
that maximize the probability of observed data).

Chebyshev's work unlocked the modern theory of limit theorems and gave
science its standard error bars. Yet foundational questions remain open. The classical
\emph{moment problem} (studied by Thomas Stieltjes and Hans Hamburger) asks
whether knowing every moment $m_1, m_2, m_3, \dots$ uniquely pins down a single
distribution. While true for distributions bounded on a finite interval, heavy-tailed
distributions on the whole line (such as the lognormal distribution) can share
the exact same infinite sequence of moments while having different probability
densities. Today, in modern machine learning and network science, many real-world
processes exhibit power-law tails whose variances or higher moments are infinite,
demanding new mathematical tools that go beyond classical moments.
\end{historybox}

\section{A Preview of Jensen's Inequality}
\label{ch09:jensen}

We have leaned repeatedly on the fact that squaring is \emph{convex} --- that
it bends upward. The general principle behind this fact explains directly
why $\E[X^2] \ge (\E[X])^2$. We will prove it in detail in the chapter on
inequalities; here is the core geometric statement.

A function $g$ is \emph{convex} if its graph lies below every chord and above
every tangent line --- it curves upward like a bowl, as $g(x) = x^2$ does.
Convexity has an immediate probabilistic consequence.

\begin{theorem}[Jensen's inequality, preview]
\label{ch09:thm-jensen}
If $g$ is a convex function and $X$ a random variable with finite mean, then
\[
  g\big(\E[X]\big) \;\le\; \E\big[g(X)\big].
\]
If $g$ is concave, the inequality reverses. Equality holds (for strictly
convex $g$) only when $X$ is a constant.
\end{theorem}

\begin{whybox}
Why is the function of the average no larger than the average of the function,
for a bowl-shaped $g$? Picture the points $(X, g(X))$ scattered along the
upward-curving graph. Averaging the $X$-values lands you at
$\big(\E[X], g(\E[X])\big)$ --- a point \emph{on} the curve. But averaging the
\emph{heights} $g(X)$ lands you somewhere among the points, and because the
curve bows downward away from any chord, the cloud of points sits \emph{above}
the curve between them. So the averaged height $\E[g(X)]$ is pulled up, above
the on-curve height $g(\E[X])$. Convexity systematically rewards spreading
out: the more $X$ varies, the higher $\E[g(X)]$ climbs above $g(\E[X])$.

A second, beautifully simple way to see this uses tangent lines. Because a convex
curve stays entirely above its tangent line drawn at $\mu = \E[X]$, we have
$g(x) \ge g(\mu) + g'(\mu)(x-\mu)$ for all $x$. Plug in the random variable $X$
and take expectations of both sides: by linearity, the linear term
$g'(\mu)\E[X-\mu]$ vanishes to zero, leaving $\E[g(X)] \ge g(\mu) = g(\E[X])$.
The curve out-climbs its tangent.
\end{whybox}

Take $g(x) = x^2$, which is convex. Jensen says
$(\E[X])^2 \le \E[X^2]$, and the gap between the two sides is precisely
$\Var(X) \ge 0$. So the non-negativity of variance is just Jensen's
inequality for the squaring function, and variance \emph{measures the size of
the Jensen gap}. That is a satisfying unification: the fact that spread is
non-negative and the fact that bowls bend upward are the same fact.

\begin{figure}[H]
\centering
\begin{tikzpicture}
\begin{axis}[
  width=10cm, height=5.6cm,
  domain=-0.4:3.4, samples=80,
  title style={font=\small},
  title={Jensen's inequality for $g(x)=x^2$},
  xlabel={$x$}, ylabel={$g(x)=x^2$},
  ymin=-0.5, ymax=10,
  grid=both, grid style={gray!15},
  legend pos=north west,
]
\addplot[sourceblue, very thick] {x^2};
\addlegendentry{$g(x)=x^2$}
% chord between x=0 and x=3
\addplot[densityred, thick, dashed] coordinates {(0,0) (3,9)};
\addlegendentry{chord}
% points: g(E[X]) on curve, E[g(X)] on chord, at midpoint 1.5
\addplot[only marks, mark=*, black] coordinates {(1.5,2.25)};
\node[anchor=north] at (axis cs:1.5,2.25) {\small $g(\E X)$};
\addplot[only marks, mark=square*, densityred] coordinates {(1.5,4.5)};
\node[anchor=south east, fill=white, inner sep=1pt] at (axis cs:1.42,4.6) {\small $\E[g(X)]$};
\draw[<->, thick, gray] (axis cs:1.7,2.25) -- (axis cs:1.7,4.5)
  node[midway,right] {\small $\Var(X)$};
\end{axis}
\end{tikzpicture}
\caption{For $X$ equal to $0$ or $3$ with equal probability, $\E[X]=1.5$. The
on-curve height $g(\E[X])=2.25$ sits below the chord height
$\E[g(X)]=4.5$. The gap, $2.25$, is exactly $\Var(X)=p(1-p)\cdot 3^2$ with
$p=1/2$. Jensen's gap is the variance.}
\label{ch09:fig-jensen}
\end{figure}

\section{Why Variance Is Everywhere}
\label{ch09:everywhere}

A short tour shows how widely variance runs through quantitative work. None of
the following needs more than the ideas of this chapter, yet together they
explain why ``minimize the variance'' is practically a mantra across the data
sciences.

\begin{itemize}
  \item \textbf{Least squares.} Linear regression fits a model by
        minimizing the sum of \emph{squared} residuals --- that is, by
        minimizing a sample variance. The choice of squares (rather than absolute
        errors) traces directly back to the same smoothness and additivity
        that make variance the standard measure of spread.
  \item \textbf{The bias--variance tradeoff.} A model's expected prediction
        error decomposes into bias squared plus variance plus irreducible
        noise. Underfitting means high bias; overfitting means high variance.
        Every learning algorithm is, at heart, negotiating this tradeoff, and
        the variance term is the one this chapter taught you to read.
  \item \textbf{Error bars and confidence.} Every reported uncertainty ---
        the $\pm$ on a measurement, the margin of error in a poll --- represents a
        standard deviation (or a multiple of one), almost always
        scaled by the $\sigma/\sqrt n$ law of Section~\ref{ch09:sumindep}.
  \item \textbf{Signal processing.} The ``power'' of a noise source is its
        variance; the signal-to-noise ratio compares the variance of the
        signal to the variance of the noise. Filtering is the art of removing
        variance you do not want.
  \item \textbf{Analysis of variance (ANOVA).} An entire branch of statistics
        is built on partitioning the total variance of a dataset into
        components attributable to different causes --- a direct application
        of the variance-of-a-sum decomposition.
\end{itemize}

\begin{aestheticbox}
If expectation answers ``what will happen on average,'' variance answers ``how
much should I trust that average.'' The first number without the second is a
forecast with no confidence attached --- a point with no error bar, a model
that says nothing about its own uncertainty. Mature quantitative reasoning
always carries both. Center and spread together provide the essential
description of an uncertain quantity.
\end{aestheticbox}

\section*{Summary}
\addcontentsline{toc}{section}{Summary}

\begin{itemize}
  \item \textbf{Variance} $\Var(X) = \E[(X-\mu)^2]$ is the average squared
        distance from the mean; \textbf{standard deviation}
        $\SD(X)=\sqrt{\Var(X)}$ returns it to the original units and is the
        natural ``unit of surprise.''
  \item \textbf{Computational formula:}
        $\Var(X) = \E[X^2] - (\E[X])^2$. Variance is never negative; it is
        zero only for a constant.
  \item \textbf{Affine rule:} $\Var(aX+b) = a^2\Var(X)$. Shifting does not
        change spread; scaling by $a$ scales the standard deviation by
        $\abs a$.
  \item \textbf{Standardization:} $Z = (X-\mu)/\sigma$ has mean $0$ and
        variance $1$; z-scores compare across scales.
  \item \textbf{Chebyshev:} for any distribution,
        $\P(\abs{X-\mu}\ge k\sigma)\le 1/k^2$ --- variance alone forces most
        mass to stay within a few standard deviations of the mean.
  \item \textbf{Sums:} $\Var(X+Y) = \Var(X)+\Var(Y)+2\Cov(X,Y)$. For
        independent (or uncorrelated) variables the covariance vanishes and
        variances add, giving the $\sigma/\sqrt n$ law for averages.
  \item \textbf{Higher moments:} skewness (third) measures asymmetry,
        kurtosis (fourth) measures tail weight.
  \item \textbf{Jensen:} $\E[X^2]\ge(\E[X])^2$ is the convexity of squaring,
        and the gap is exactly the variance.
\end{itemize}

\subsection*{Looking Ahead: Crossing into the Continuum}

Take a breath and look back across the ground we have covered. In the last two
chapters, you established the two master coordinates of probability:
expectation (the center of mass) and variance (the moment of inertia). But notice
the quiet assumption running beneath every computation we performed: our random
variables lived on tidy, discrete lists of values. We rolled dice, flipped coins,
and tallied heads. Every expectation was a sum $\mu = \sum x\,\P(X=x)$, and every
variance was a sum of squared deviations $\sigma^2 = \sum (x-\mu)^2\,\P(X=x)$.

The physical world, however, rarely stands on discrete steps. At what exact
microsecond will a radioactive particle decay? What is the precise angle of a
spinning wheel when it halts? What is the exact temperature of a cup of coffee
as it cools? These quantities do not jump from one integer to the next; they
vary continuously across an unbroken continuum of real numbers.

In Chapter~10, we move from discrete variables to continuous ones. For a continuous
variable, the probability of landing on any single exact number is zero ($\P(X=x) = 0$).
To work with outcomes on an interval, probability mass becomes a continuous density
function, and sums $\sum$ are replaced by integrals $\int$.

The core ideas carry over: expectation is still the center of mass $\int x\,f(x)\,\dd x$,
variance is still $\int (x-\mu)^2 f(x)\,\dd x$, and the standard deviation still measures
distance from the mean.

\section{Exercises}
\label{ch09:exercises}

\begin{exercise}[Warm-up computations]
\label{ch09:exr-warmup}
For each of the following, compute $\E[X]$, $\E[X^2]$, $\Var(X)$, and
$\SD(X)$.
\begin{enumerate}
  \item $X$ takes values $2, 4, 6$ each with probability $1/3$.
  \item $X \sim \Bern(0.2)$.
  \item $X$ is the number of heads in two independent fair coin flips.
\end{enumerate}
\end{exercise}

\begin{exercise}[Shifting and scaling]
\label{ch09:exr-affine}
A temperature $C$ measured in Celsius has mean $20$ and standard deviation
$5$. The Fahrenheit reading is $F = \tfrac{9}{5}C + 32$. Find $\E[F]$,
$\Var(F)$, and $\SD(F)$. Explain in words why the $+32$ did not affect the
spread.
\end{exercise}

\begin{exercise}[The computational formula in reverse]
\label{ch09:exr-reverse}
A random variable has mean $\E[X] = 3$ and variance $\Var(X) = 4$. Find
$\E[X^2]$. Then find $\E[(2X-1)^2]$. (Hint: expand and reuse the first two
moments.)
\end{exercise}

\begin{exercise}[z-scores in practice]
\label{ch09:exr-zscore}
A standardized test has mean $500$ and standard deviation $100$; a second test
has mean $24$ and standard deviation $4$. Alice scores $640$ on the first;
Bob scores $30$ on the second. Convert both to z-scores and say who performed
better relative to peers. By how many points above the mean did each score,
in units of their own test's standard deviation?
\end{exercise}

\begin{exercise}[Variance of a sum, with and without independence]
\label{ch09:exr-sum}
Let $X$ and $Y$ each have variance $9$.
\begin{enumerate}
  \item If $X$ and $Y$ are independent, find $\Var(X+Y)$ and $\SD(X+Y)$.
  \item If instead $\Cov(X,Y) = 6$, find $\Var(X+Y)$.
  \item If $Y = X$ (perfect dependence), find $\Var(X+Y)$ directly and
        confirm it matches the formula with the appropriate covariance.
\end{enumerate}
\end{exercise}

\begin{exercise}[The $\sqrt n$ law]
\label{ch09:exr-sqrtn}
A single measurement of a length has standard deviation $\sigma = 2$ mm.
\begin{enumerate}
  \item How many independent measurements must you average so that the
        standard deviation of the average is at most $0.25$ mm?
  \item Suppose you have reached an average error of $0.5$ mm with some number
        of measurements. How many \emph{times} more measurements are needed to
        reach $0.1$ mm?
\end{enumerate}
\end{exercise}

\begin{exercise}[Chebyshev in action]
\label{ch09:exr-cheby}
A machine fills bags with a mean weight of $\mu = 500$ g and a standard
deviation of $\sigma = 4$ g; the distribution of weights is unknown.
\begin{enumerate}
  \item Use Chebyshev's inequality to give a guaranteed lower bound on the
        probability that a bag weighs between $492$ g and $508$ g.
  \item What is the largest possible probability that a bag deviates from
        $500$ g by $12$ g or more?
  \item If the weights were known to be normal, roughly what fraction would
        fall within $492$--$508$ g? Compare to your Chebyshev bound and explain
        the gap.
\end{enumerate}
\end{exercise}

\begin{exercise}[Binomial variance two ways]
\label{ch09:exr-binom}
Let $X \sim \Bin(n,p)$.
\begin{enumerate}
  \item Using the decomposition $X = \sum_{i=1}^n B_i$ into independent
        Bernoulli trials, state $\Var(X)$.
  \item For $n = 100$, find the value of $p$ that maximizes $\Var(X)$, and
        give the resulting standard deviation.
  \item Interpret: for which $p$ is a binomial outcome hardest to predict?
\end{enumerate}
\end{exercise}

\begin{exercise}[Skewness of a simple variable]
\label{ch09:exr-skew}
Let $X$ take the value $0$ with probability $0.9$ and the value $10$ with
probability $0.1$.
\begin{enumerate}
  \item Compute $\mu = \E[X]$ and $\sigma = \SD(X)$.
  \item Compute the third central moment $\E[(X-\mu)^3]$ and the skewness.
  \item Is the distribution skewed left or right? Sketch why this matches the
        sign you found.
\end{enumerate}
\end{exercise}

\begin{exercise}[Jensen and the variance gap]
\label{ch09:exr-jensen}
Let $X$ take values $1$ and $9$ with equal probability.
\begin{enumerate}
  \item Compute $g(\E[X])$ and $\E[g(X)]$ for $g(x) = x^2$, and verify that
        their difference equals $\Var(X)$.
  \item The square-root function $h(x) = \sqrt x$ is concave. Compute
        $h(\E[X])$ and $\E[h(X)]$ and check that the inequality runs the
        other way, as Jensen predicts for concave functions.
\end{enumerate}
\end{exercise}

\begin{exercise}[Proof: variance is translation invariant and quadratic in
scale]
\label{ch09:exr-proof}
Starting from the definition $\Var(X) = \E[(X-\E[X])^2]$ and using only
linearity of expectation, prove that $\Var(aX+b) = a^2\Var(X)$ for all
constants $a,b$. Then prove the computational formula
$\Var(X) = \E[X^2]-(\E[X])^2$ as a special consequence of your expansion.
\end{exercise}

\begin{exercise}[Proof: the variance of a sum]
\label{ch09:exr-proofsum}
Let $X_1,\dots,X_n$ be random variables with finite variances. Prove that
\[
  \Var\!\Big(\sum_{i=1}^n X_i\Big)
  = \sum_{i=1}^n \Var(X_i) + \sum_{i\ne j}\Cov(X_i,X_j).
\]
Deduce the special case for pairwise uncorrelated variables. Where exactly in
your proof did you use (or not use) independence?
\end{exercise}

\begin{exercise}[Conceptual: when does averaging \emph{not} help?]
\label{ch09:exr-concept}
The $\sigma/\sqrt n$ law assumed the measurements were independent. Describe a
realistic scenario --- in measurement, polling, or finance --- where averaging
many observations fails to reduce the error nearly as fast as $1/\sqrt n$.
Identify which term in the variance-of-a-sum formula is responsible, and state
what its sign must be for the failure to occur.
\end{exercise}

\begin{exercise}[Mental check: impossible moments]
\label{ch09:exr-impossible}
A student claims to have analyzed a dataset and found a random variable with
mean $\E[X] = 10$ and second raw moment $\E[X^2] = 95$.
Explain in a single sentence why this claim is mathematically impossible,
without needing to know anything else about the distribution.
\end{exercise}

\begin{exercise}[Chebyshev's leash is sharp]
\label{ch09:exr-cheby-sharp}
Chebyshev's inequality guarantees that $\P(\abs{X-\mu}\ge k\sigma) \le 1/k^2$.
Show that this bound cannot be made any tighter without additional assumptions,
by constructing a distribution that hits the bound with exact equality.
Specifically, for $k=2$, define a discrete variable $X$ that takes values:
\[
  X = \begin{cases}
    -2 & \text{with probability } \tfrac{1}{8}, \\
    0  & \text{with probability } \tfrac{3}{4}, \\
    +2 & \text{with probability } \tfrac{1}{8}.
  \end{cases}
\]
\begin{enumerate}
  \item Compute $\E[X]$ and show that $\sigma = \SD(X) = 1$.
  \item Calculate the exact probability $\P(\abs{X-\mu} \ge 2\sigma)$, and confirm
        that it equals $1/k^2 = 1/4$ exactly.
\end{enumerate}
\end{exercise}
